% Which is more accurate, LIGHT v0, v0_1 or v0_2. The three versions on the
% same tests. No number here is typed by hand. They come from seven
% generated files: v0's map and overview numbers
% (docs/system_map_basic_numbers.tex, docs/overview_fast.tex,
% docs/overview_slow.tex), v0_1's (docs/system_map_v0_1_numbers.tex,
% docs/overview_v0_1_fast.tex, docs/overview_v0_1_slow.tex) and this
% comparison's review_numbers.tex (make_review.py), which also holds v0_2's
% column. Build from LIGHT/docs/REVIEW:
%   python3 make_review.py && latexmk -pdf v0_vs_v0_1_vs_v0_2.tex
\documentclass[11pt]{article}
\usepackage[a4paper,margin=2.2cm]{geometry}
\usepackage[T1]{fontenc}
\usepackage[scaled=0.95]{helvet}
\renewcommand{\familydefault}{\sfdefault}
\usepackage{amsmath}
\usepackage{sansmath}
\usepackage{booktabs,tabularx,array}
\usepackage{xcolor}
\usepackage{float}
\usepackage{needspace}
\setlength{\parindent}{0pt}
\setlength{\parskip}{5pt}
\definecolor{ink2}{HTML}{52514E}
% a page of tables starts at its top
\makeatletter\setlength{\@fptop}{0pt}\makeatother

\input{../system_map_basic_numbers}
\input{../overview_fast}
\input{../overview_slow}
\input{../system_map_v0_1_numbers}
\input{../overview_v0_1_fast}
\input{../overview_v0_1_slow}
% LIGHT v0, v0_1 and v0_2 comparison numbers. GENERATED by make_review.py
% from v0 (0.1.0), v0_1 (0.1.1) and v0_2 (0.1.2). Do not edit. Rerun the script instead.
\providecommand\rvdef[2]{\expandafter\def\csname rv@#1\endcsname{#2}}
\rvdef{v0.version}{0.1.0}
\rvdef{v01.version}{0.1.1}
\rvdef{v02.version}{0.1.2}
\rvdef{tah.v0.def.g0}{2.05}
\rvdef{tah.v0.def.x}{0.32}
\rvdef{tah.v0.def.te}{1.5}
\rvdef{tah.v0.def.tg}{372}
\rvdef{tah.v0.defkr2.g0}{2.04}
\rvdef{tah.v0.defkr2.x}{0.32}
\rvdef{tah.v0.defkr2.te}{1.5}
\rvdef{tah.v0.defkr2.tg}{372}
\rvdef{tah.v0.own.g0}{6.60}
\rvdef{tah.v0.own.x}{1.03}
\rvdef{tah.v0.own.te}{3.8}
\rvdef{tah.v0.own.tg}{431}
\rvdef{tah.v0.noF.g0}{5.71}
\rvdef{tah.v0.noF.x}{0.89}
\rvdef{tah.v0.noF.te}{1.6}
\rvdef{tah.v0.noF.tg}{449}
\rvdef{tah.v01.def.g0}{5.94}
\rvdef{tah.v01.def.x}{0.93}
\rvdef{tah.v01.def.te}{1.6}
\rvdef{tah.v01.def.tg}{447}
\rvdef{tah.v01.defkr2.g0}{6.10}
\rvdef{tah.v01.defkr2.x}{0.95}
\rvdef{tah.v01.defkr2.te}{1.6}
\rvdef{tah.v01.defkr2.tg}{448}
\rvdef{tah.v01.own.g0}{6.56}
\rvdef{tah.v01.own.x}{1.03}
\rvdef{tah.v01.own.te}{3.8}
\rvdef{tah.v01.own.tg}{428}
\rvdef{tah.v01.noF.g0}{5.69}
\rvdef{tah.v01.noF.x}{0.89}
\rvdef{tah.v01.noF.te}{1.6}
\rvdef{tah.v01.noF.tg}{446}
\rvdef{tah.v02.def.g0}{5.69}
\rvdef{tah.v02.def.x}{0.89}
\rvdef{tah.v02.def.te}{1.6}
\rvdef{tah.v02.def.tg}{446}
\rvdef{tah.v02.own.g0}{5.53}
\rvdef{tah.v02.own.x}{0.87}
\rvdef{tah.v02.own.te}{3.8}
\rvdef{tah.v02.own.tg}{423}
\rvdef{tah.v02.kannari.g0}{6.10}
\rvdef{tah.v02.kannari.x}{0.95}
\rvdef{tah.v02.kannari.te}{1.6}
\rvdef{tah.v02.kannari.tg}{448}
\rvdef{tah.v02.hiF.g0}{6.23}
\rvdef{tah.v02.hiF.x}{0.97}
\rvdef{tah.v02.hiF.te}{1.6}
\rvdef{tah.v02.hiF.tg}{449}
\rvdef{tah.v02.tri.g0}{2.84}
\rvdef{tah.v02.tri.x}{0.44}
\rvdef{tah.v02.tri.te}{1.6}
\rvdef{tah.v02.tri.tg}{447}
\rvdef{v02.xu.50}{0.91}
\rvdef{v02.xu.100}{0.97}
\rvdef{v02.sc.eden_kr}{2.3}
\rvdef{v02.sc.eden_f2}{3.2}
\rvdef{v02.sc.ksu_fig2}{4.2}
\rvdef{v02.sc.ksu_fig1}{6.1}
\rvdef{v02.sc.xu_figs5_7}{11.7}
\rvdef{v02.sc.hakuta_fig3}{34.1}
\rvdef{v02.sc.hakuta_fig2}{19.3}
\rvdef{v02.sc.ys_ir300}{\ensuremath{-}0.5}
\rvdef{v02.sc.ys_ratio_230_300}{\ensuremath{-}8.4}
\rvdef{v02.sc.bs_tau_b}{\ensuremath{-}22.3}
\rvdef{v02.tah.g0}{5.69}
\rvdef{v02.tah.g0.x}{0.89}
\rvdef{v02.tah.isat}{1.55}
\rvdef{v02.tah.isat.x}{1.11}
\rvdef{v02.tah.eta}{0.268}
\rvdef{v02.tah.eta.x}{0.99}
\rvdef{v02.tah.tg}{446}
\rvdef{v02.pure.Kr.x}{1.02}
\rvdef{v02.pure.Ar.x}{1.03}
\rvdef{v02.pure.dr.lo}{1.00}
\rvdef{v02.pure.dr.hi}{1.00}
\rvdef{v02.pure.k15}{$2.4{\times}10^{-31}$}
\rvdef{v02.pure.vuv}{0.57}
\rvdef{v02.post.g0.lo}{4.99}
\rvdef{v02.post.g0.mid}{5.39}
\rvdef{v02.post.g0.hi}{6.18}
\rvdef{v02.el.dev.te}{0.7}
\rvdef{v02.el.dev.ne}{2.5}
\rvdef{ref.v0}{5.24}
\rvdef{ref.v0.tg}{604}
\rvdef{ref.v0.vuv}{0.0}
\rvdef{ref.v0groups}{4.85}
\rvdef{ref.v0groups.tg}{762}
\rvdef{ref.v0groups.vuv}{0.0}
\rvdef{ref.v01}{4.73}
\rvdef{ref.v01.tg}{762}
\rvdef{ref.v01.vuv}{0.0}
\rvdef{ref.noBA}{4.90}
\rvdef{ref.noBA.tg}{762}
\rvdef{ref.noBA.vuv}{0.0}
\rvdef{ref.noFF}{4.69}
\rvdef{ref.noFF.tg}{762}
\rvdef{ref.noFF.vuv}{0.0}
\rvdef{ref.v0ions}{4.75}
\rvdef{ref.v0ions.tg}{762}
\rvdef{ref.v0ions.vuv}{0.0}
\rvdef{ref.noIIT}{4.73}
\rvdef{ref.noIIT.tg}{763}
\rvdef{ref.noIIT.vuv}{0.0}
\rvdef{ref.v0E}{4.73}
\rvdef{ref.v0E.tg}{763}
\rvdef{ref.v0E.vuv}{0.0}
\rvdef{ref.v02}{4.67}
\rvdef{ref.v02.tg}{649}
\rvdef{ref.v02.vuv}{20.0}
\rvdef{ref.v02kannari}{9.17}
\rvdef{ref.v02kannari.tg}{724}
\rvdef{ref.v02kannari.vuv}{1.3}
\rvdef{ref.v02pen}{4.65}
\rvdef{ref.v02pen.tg}{704}
\rvdef{ref.v02pen.vuv}{10.8}
\rvdef{ref.gap}{0.51}
\rvdef{ref.d.network}{0.39}
\rvdef{ref.d.values}{0.12}
\rvdef{ref.gap12}{0.07}
\rvdef{ref.cool12}{113}
\rvdef{band.kr2_f.lo}{4.67}
\rvdef{band.kr2_f.hi}{9.50}
\rvdef{band.kr2_pen.lo}{4.67}
\rvdef{band.kr2_pen.hi}{4.65}
\rvdef{band.kr2f_f.lo}{4.67}
\rvdef{band.kr2f_f.hi}{4.87}
\rvdef{band.kr2_form.lo}{4.67}
\rvdef{band.kr2_form.hi}{4.52}
\rvdef{band.kr2_f2.lo}{4.67}
\rvdef{band.kr2_f2.hi}{4.60}
\rvdef{band.kr_pen.lo}{4.87}
\rvdef{band.kr_pen.hi}{4.67}
\rvdef{band.kr2f_f2.lo}{4.67}
\rvdef{band.kr2f_f2.hi}{4.88}
\rvdef{band.trimer_tg_pow.lo}{4.67}
\rvdef{band.trimer_tg_pow.hi}{2.22}
\rvdef{band.minor}{0.21}
\rvdef{band.n}{8}
\rvdef{grid.diff}{91}
\rvdef{grid.diff12}{11}
\rvdef{grid.spread}{91}
\rvdef{grid.kr2}{2.7}
\rvdef{grid.tri}{66}
\rvdef{grid.n}{15}
\rvdef{grid.nband}{10}
\rvdef{grid.hotmin}{1173}
\def\rvBands{\Krts{} + F \ra KrF(B) + Kr & 0 & 0 to $1.0{\times}10^{-9}$\unit{cm^3\,s^{-1}} & 4.67 to 9.50 & 20.0 to 0.5 & 649 to 726\\
\Krts{} + \Krts{} \ra \Krtp{} + 2Kr + \el{} & 0 & 0 to $5.0{\times}10^{-10}$\unit{cm^3\,s^{-1}} & 4.67 to 4.65 & 20.0 to 10.8 & 649 to 704\\
\KtF{} + F \ra 2Kr + 2F & 0 & 0 to $1.8{\times}10^{-10}$\unit{cm^3\,s^{-1}} & 4.67 to 4.87 & 20.0 to 19.9 & 649 to 670\\
\Krs{} + 2Kr \ra \Krts{} + Kr, 300\unit{K} & $1.7{\times}10^{-32}$ & $1.7{\times}10^{-32}$ to $5.5{\times}10^{-32}$\unit{cm^6\,s^{-1}} & 4.67 to 4.52 & 20.0 to 28.0 & 649 to 601\\
\Krts{} + \Ft{}, both channels & $2.1{\times}10^{-10}$ & $2.1{\times}10^{-10}$ to $5.0{\times}10^{-10}$\unit{cm^3\,s^{-1}} & 4.67 to 4.60 & 20.0 to 19.1 & 649 to 654\\
\Krs{} + \Krs{} \ra \Krp{} + Kr + \el{} & $5.0{\times}10^{-10}$ & $1.5{\times}10^{-10}$ to $5.0{\times}10^{-10}$\unit{cm^3\,s^{-1}} & 4.87 to 4.67 & 25.0 to 20.0 & 616 to 649\\
\KtF{} + \Ft{} & $1.2{\times}10^{-10}$ & $1.2{\times}10^{-10}$ to $4.3{\times}10^{-10}$\unit{cm^3\,s^{-1}} & 4.67 to 4.88 & 20.0 to 19.9 & 649 to 663\\
the trimer's $(T_g/300)^n$ & $-3.32$ & $-3.32$ to $-0.50$ & 4.67 to 2.22 & 20.0 to 19.5 & 649 to 658\\}
\def\rvGrid{0.1 & 100 & 3.73 & 3.46 & 3.47 & 3.55 & 3.39 & 1173\\
0.1 & 300 & 1.25 & 0.80 & 0.80 & 0.81 & 0.79 & 2434\\
0.1 & 500 & 0.75 & 0.39 & 0.39 & 0.40 & 0.39 & 3458\\
0.25 & 100 & 8.43 & 9.22 & 9.23 & 10.33 & 8.07 & 826\\
0.25 & 300 & 2.83 & 2.39 & 2.40 & 2.61 & 2.15 & 2044\\
0.25 & 500 & 1.71 & 1.21 & 1.21 & 1.30 & 1.10 & 3102\\
0.5 & 100 & 11.31 & 14.21 & 13.84 & 16.13 & 9.25 & 532\\
0.5 & 300 & 4.38 & 4.11 & 4.13 & 5.86 & 2.93 & 1191\\
0.5 & 500 & 2.63 & 2.24 & 2.26 & 3.03 & 1.65 & 1928\\
1 & 100 & 4.79 & 5.73 & 5.42 & 5.84 & 3.28 & 440\\
1 & 300 & 5.24 & 4.73 & 4.67 & 9.50 & 2.22 & 649\\
1 & 500 & 3.09 & 2.72 & 2.71 & 7.18 & 1.33 & 915\\
2 & 100 & 1.00 & 1.25 & 1.17 & 1.28 & 0.83 & 378\\
2 & 300 & 2.30 & 2.45 & 2.17 & 2.64 & 0.89 & 519\\
2 & 500 & 2.81 & 2.31 & 2.12 & 4.48 & 0.73 & 604\\}

\makeatletter
\newcommand\lm[1]{\ifcsname lm@#1\endcsname\csname lm@#1\endcsname
  \else\errmessage{no v0 map number '#1'}\fi}
\newcommand\ov[1]{\ifcsname ov@#1\endcsname\csname ov@#1\endcsname
  \else\errmessage{no v0 overview number '#1'}\fi}
\newcommand\vm[1]{\ifcsname vm@#1\endcsname\csname vm@#1\endcsname
  \else\errmessage{no v0_1 map number '#1'}\fi}
\newcommand\vi[1]{\ifcsname vi@#1\endcsname\csname vi@#1\endcsname
  \else\errmessage{no v0_1 overview number '#1'}\fi}
\newcommand\rv[1]{\ifcsname rv@#1\endcsname\csname rv@#1\endcsname
  \else\errmessage{no review number '#1'. Rerun make_review.py}\fi}
\makeatother

% chemistry and units
\newcommand{\unit}[1]{\ensuremath{\,\mathrm{#1}}}
\newcommand{\el}{\ensuremath{\mathrm{e}^-}}
\newcommand{\Krs}{Kr\ensuremath{^*}}
\newcommand{\Krp}{Kr\ensuremath{^+}}
\newcommand{\Krtp}{Kr\ensuremath{_2^+}}
\newcommand{\Krts}{Kr\ensuremath{_2^*}}
\newcommand{\KtF}{Kr\ensuremath{_2}F\ensuremath{^*}}
\newcommand{\Ft}{F\ensuremath{_2}}
\newcommand{\ra}{\ensuremath{\to}\ }
\newcommand{\Edep}{\ensuremath{E_\mathrm{dep}}}
\newcommand{\gIP}{\ensuremath{g_0I_\mathrm{sat}/P_\mathrm{in}}}
\newcommand{\x}[1]{\ensuremath{\times}#1}

\begin{document}
\sansmath

{\LARGE\bfseries Which is more accurate, LIGHT v0, v0\_1 or v0\_2}\par
\vspace{4pt}
{\large The three versions on the same tests. v0 is version \lm{version},
v0\_1 is version \vm{version} and v0\_2 is version \rv{v02.version}, 29
September 2026.}\par
\vspace{10pt}

\noindent\fbox{\parbox{\dimexpr\linewidth-2\fboxsep-2\fboxrule\relax}{%
\setlength{\parskip}{5pt}%
\textbf{Bottom line.} v0\_2 is the best of the three models. It is more
accurate than v0, it scores level with v0\_1, and it is the only version
whose default network rests on measurements and theory alone.

v0 falls behind because its default network holds $n_e$ fixed and has no
charge chemistry, and under Tahoma's beam it gives the gain at
\x{\rv{tah.v0.def.x}} the measured value. v0\_1 and v0\_2 give the same
numbers on every kinetic dataset, every pure-gas check and the electron
closure. At Tahoma they trade places. v0\_2 gives the formation efficiency
at \x{\rv{v02.tah.eta.x}}, the nearest of the three, and the gain at
\x{\rv{tah.v02.def.x}}. v0\_1 comes closer on the gain,
\x{\rv{tah.v01.def.x}} by default and \x{\vi{tah.g0.x}} as ML scores it,
but only by leaving measurements behind. Its default network has no
\Krts{} route and so leaves out \Krs{} + 2Kr, a measured channel that takes
\Krs{} from the harpoon reaction. ML's configuration adds the route with
Kannari's estimates for \Krts{} + F and \Krts{} + \Krts{}, two rates that
have never been measured, and v0\_2 with those two estimates gives v0\_1's
\rv{tah.v02.kannari.g0}\unit{\%/cm}. Under the predictive-only policy an
estimate may not set a rate. v0\_2 also gives every unmeasured or disputed
rate as a band, so it says how far its answer can move, up to a factor
\rv{grid.kr2} through \Krts{} + F and down by \rv{grid.tri}\,\% through the
trimer's temperature law across the notes' ranges.

Most of the gain v0\_2 still misses at Tahoma lies inside its \Krts{} + F
band, whose top gives \x{\rv{tah.v02.hiF.x}}. A measurement of that rate
would show whether the gap is that reaction.}}
\medskip

v0\_1 kept v0's KrF(B,C) block and every formation link, and v0\_2 kept
v0\_1's, so wherever they run the same chemistry they score alike, and
Table~\ref{tab:same} lists every test. On the 300\unit{K} kinetic datasets
the three agree to the digit. Under Tahoma's electron beam, with the ion
and F$_2$-electron chemistry on, $n_e$ free and the \Krts{} group with
Kannari's estimates in it, v0 gives the small-signal gain at
\x{\ov{tah.g0.x}} the measured \ov{tah.g0.meas}\unit{\%/cm} and v0\_1 at
\x{\vi{tah.g0.x}}. Yu and Setser's C/B ratio between 230 and 300\unit{K}
improves from \ov{sc.ys_ratio_230_300} to \vi{sc.ys_ratio_230_300}\,\%,
but v0\_1 and v0\_2 book their B-C gap, which rests partly on that ratio,
so the improvement is not an independent test.

The difference between v0 and the others is what each runs by default.
v0's basic network holds $n_e$ at $10^{15}$\unit{cm^{-3}} and has no charge
chemistry, so every ion it makes stays an ion. Under Tahoma's beam, which
makes ion pairs, that network gives the gain at \x{\rv{tah.v0.def.x}} the
measured value, and \x{\rv{tah.v0.defkr2.x}} with the \Krts{} group added.
v0's overview scored \x{\ov{tah.g0.x}} because ML ran v0 with its optional
ion, F$_2$-electron and \Krts{} groups on and $n_e$ free, and the overview
did not say so. v0\_1 made that chemistry its default, and its default
network gives \x{\rv{tah.v01.def.x}}. v0\_2's gives \x{\rv{tah.v02.def.x}}.

Between v0\_1 and v0\_2 the difference at Tahoma is the \Krts{} route.
v0\_2 carries its measured rows, Ueda's \Krs{} + 2Kr, the 270-ns lifetime
and Quigley and Hughes' \Krts{} + F$_2$. \Krts{} + F \ra KrF(B) + Kr and
\Krts{} + \Krts{} Penning ionization have no measurement, so v0\_2 leaves
them at zero and gives the gain at \x{\rv{tah.v02.def.x}}. v0\_1 as ML
scores it, with Kannari's estimates for those two, gives
\x{\vi{tah.g0.x}}, and v0\_2 with the same two estimates gives
\x{\rv{tah.v02.kannari.x}}, so those two rates are the whole difference.
On the formation efficiency the order turns, \x{\rv{v02.tah.eta.x}} for
v0\_2 against \x{\vi{tah.eta.x}} for v0\_1, and the saturation
intensities agree, \vi{tah.isat} and \rv{v02.tah.isat}\unit{MW/cm^2}. Over the posterior of the B/C constants v0\_2's gain
spans \rv{v02.post.g0.lo} to \rv{v02.post.g0.hi}\unit{\%/cm} and v0\_1's
\vi{post.g0.lo} to \vi{post.g0.hi}. Tahoma's gain would favour a
nonzero \Krts{} + F and its formation efficiency the zero, but Tahoma is
device data, so it scores that rate and never sets it. At the top of its
band v0\_2 gives the gain at \x{\rv{tah.v02.hiF.x}}.

Two unmeasured quantities dominate what all three versions leave open. One is
the trimer's temperature law above \vi{tri.klimek.hi}\unit{K}. At Tahoma's
\rv{v02.tah.tg}\unit{K} the weak $T_g^{-1/2}$ law gives v0\_2's gain at
\x{\rv{tah.v02.tri.x}} and v0\_1's at \x{\vi{tah.tri.x}}, so Tahoma's
gain counts against that law, as a score, and at the reference point it
takes v0\_2's 248-nm share from \rv{band.trimer_tg_pow.lo} to
\rv{band.trimer_tg_pow.hi}\,\%. The other is \Krts{} + F, whose band
takes the same share from \rv{band.kr2_f.lo} to \rv{band.kr2_f.hi}\,\%.
Each of v0\_2's other bands moves it by \rv{band.minor} points or less. v0\_2's changes are in
\texttt{LIGHT/v0\_2/README.md}, and v0\_1's accuracy is written up in
\texttt{docs/accuracy\_overview\_v0\_1.pdf}.

\begin{table}[htbp]
\centering\footnotesize
\renewcommand{\arraystretch}{1.12}
\begin{tabularx}{\linewidth}{@{}>{\raggedright\arraybackslash}X
  >{\raggedright\arraybackslash}p{1.9cm}>{\raggedright\arraybackslash}p{2.35cm}
  >{\raggedright\arraybackslash}p{2.35cm}
  >{\raggedright\arraybackslash}p{2.35cm}@{}}
\toprule
test & measured & v0 & v0\_1 & v0\_2\\
\midrule
KrF(B) decay against Kr and F$_2$, Eden 1978 & & rms \ov{sc.eden_kr} and
  \ov{sc.eden_f2}\,\% & rms \vi{sc.eden_kr} and \vi{sc.eden_f2}\,\% & rms
  \rv{v02.sc.eden_kr} and \rv{v02.sc.eden_f2}\,\%\\
KrF(B,C) slow decay, Gadomski 1992 & & rms \ov{sc.ksu_fig1} and
  \ov{sc.ksu_fig2}\,\% & rms \vi{sc.ksu_fig1} and \vi{sc.ksu_fig2}\,\% &
  rms \rv{v02.sc.ksu_fig1} and \rv{v02.sc.ksu_fig2}\,\%\\
C/B emission ratio at 300\unit{K}, Yu and Setser 1990 & \ov{ys.meas}
  $\pm$ \ov{ys.err} & \ov{sc.ys_ir300}\,\% & \vi{sc.ys_ir300}\,\% &
  \rv{v02.sc.ys_ir300}\,\%\\
KrF(B,C) decays, Xu 1993 Figs. 5 to 7 & & rms \ov{sc.xu_figs5_7}\,\% &
  rms \vi{sc.xu_figs5_7}\,\% & rms \rv{v02.sc.xu_figs5_7}\,\%\\
decay constants Xu 1993 printed & \ov{xu.50.meas} and
  \ov{xu.100.meas}\unit{s^{-1}} & \x{\ov{xu.50}} and \x{\ov{xu.100}} &
  \x{\vi{xu.50}} and \x{\vi{xu.100}} & \x{\rv{v02.xu.50}} and
  \x{\rv{v02.xu.100}}\\
C-A decays, Hakuta 1990 & & rms \ov{sc.hakuta_fig2} and
  \ov{sc.hakuta_fig3}\,\% & rms \vi{sc.hakuta_fig2} and
  \vi{sc.hakuta_fig3}\,\% & rms \rv{v02.sc.hakuta_fig2} and
  \rv{v02.sc.hakuta_fig3}\,\%\\
$\tau_B$, Burnham and Searles 1977 & \ov{sc.bs.meas} $\pm$
  \ov{sc.bs.err}\unit{ns} & \ov{sc.bs_tau_b}\,\% & \vi{sc.bs_tau_b}\,\% &
  \rv{v02.sc.bs_tau_b}\,\%\\
C/B ratio 230\unit{K} over 300\unit{K}, Yu and Setser & \ov{ysr.meas}
  $\pm$ \ov{ysr.err} & \ov{sc.ys_ratio_230_300}\,\% &
  \vi{sc.ys_ratio_230_300}\,\%, not independent &
  \rv{v02.sc.ys_ratio_230_300}\,\%, not independent\\
\midrule
\Krs{} + \Krp{} made per eV, pure Kr and Ar, Saito 2003 & &
  \x{\ov{pure.Kr.x}} and \x{\ov{pure.Ar.x}} & \x{\vi{pure.Kr.x}} and
  \x{\vi{pure.Ar.x}} & \x{\rv{v02.pure.Kr.x}} and \x{\rv{v02.pure.Ar.x}}\\
\Krtp{} + \el{} recombination, Shiu and Biondi 1977 & &
  \x{\ov{pure.dr.lo}} to \x{\ov{pure.dr.hi}} & \x{\vi{pure.dr.lo}}, the
  input & \x{\rv{v02.pure.dr.lo}}, the input\\
\Krp{} + 2Kr \ra \Krtp{} + Kr, Smith et al. 1972 & \ov{pure.conv.lo} to
  \ov{pure.conv.hi} & \ov{pure.k15} & \vi{pure.k15}, a mean with two
  others & \rv{v02.pure.k15}, the same mean\\
VUV yield of pure Kr at low pump, Morozov 2008 & \ov{pure.vuv.meas}
  $\pm$ \ov{pure.vuv.err} & \ov{pure.vuv} & \vi{pure.vuv} &
  \rv{v02.pure.vuv}\\
electron closure at Tahoma, $T_e$ and $n_e$ against the kinetic loop & &
  within \ov{el.dev.te} and \ov{el.dev.ne}\,\% & within \vi{el.dev.te}
  and \vi{el.dev.ne}\,\% & within \rv{v02.el.dev.te} and
  \rv{v02.el.dev.ne}\,\%\\
\midrule
Tahoma small-signal gain $g_0$, each version's default network &
  \ov{tah.g0.meas}\unit{\%/cm} & \rv{tah.v0.def.g0}, \x{\rv{tah.v0.def.x}}
  & \rv{tah.v01.def.g0}, \x{\rv{tah.v01.def.x}} & \rv{tah.v02.def.g0},
  \x{\rv{tah.v02.def.x}}\\
\quad the same with the \Krts{} group & & \rv{tah.v0.defkr2.g0},
  \x{\rv{tah.v0.defkr2.x}} & \rv{tah.v01.defkr2.g0},
  \x{\rv{tah.v01.defkr2.x}} & in the default\\
\quad as ML scores each version, with ion and F$_2$-electron chemistry,
  \Krts{} and free $n_e$ & & \ov{tah.g0}, \x{\ov{tah.g0.x}} & \vi{tah.g0},
  \x{\vi{tah.g0.x}} & \rv{v02.tah.g0}, \x{\rv{v02.tah.g0.x}}, the
  default\\
\quad the same without the estimated \Krts{} + F \ra KrF & &
  \rv{tah.v0.noF.g0}, \x{\rv{tah.v0.noF.x}} & \rv{tah.v01.noF.g0},
  \x{\rv{tah.v01.noF.x}} & as above\\
\quad with Kannari's \Krts{} + F and \Krts{} + \Krts{}, and at the top of
  the \Krts{} + F band & & & & \rv{tah.v02.kannari.g0},
  \x{\rv{tah.v02.kannari.x}}, and \rv{tah.v02.hiF.g0},
  \x{\rv{tah.v02.hiF.x}}\\
\quad over the posterior of the B/C constants & & \ov{post.g0.mid}
  [\ov{post.g0.lo}, \ov{post.g0.hi}] & \vi{post.g0.mid} [\vi{post.g0.lo},
  \vi{post.g0.hi}] & \rv{v02.post.g0.mid} [\rv{v02.post.g0.lo},
  \rv{v02.post.g0.hi}]\\
\quad on each version's own power input, no beam & &
  \rv{tah.v0.own.g0}, \x{\rv{tah.v0.own.x}} & \rv{tah.v01.own.g0},
  \x{\rv{tah.v01.own.x}} & \rv{tah.v02.own.g0}, \x{\rv{tah.v02.own.x}}\\
Tahoma saturation intensity $I_\mathrm{sat}$ & \ov{tah.isat.meas} $\pm$
  \ov{tah.isat.err}\unit{MW/cm^2} & \ov{tah.isat}, \x{\ov{tah.isat.x}} &
  \vi{tah.isat}, \x{\vi{tah.isat.x}} & \rv{v02.tah.isat},
  \x{\rv{v02.tah.isat.x}}\\
Tahoma formation efficiency \gIP & \ov{tah.eta.meas} & \ov{tah.eta},
  \x{\ov{tah.eta.x}} & \vi{tah.eta}, \x{\vi{tah.eta.x}} &
  \rv{v02.tah.eta}, \x{\rv{v02.tah.eta.x}}\\
\bottomrule
\end{tabularx}
\caption{v0, v0\_1 and v0\_2 on the same tests. A percentage is the model's
deviation from the data, signed for a single value and rms for a set. A
\x{} factor is the model over the measurement. Tahoma's gain is in
\%/cm. For v0\_1, ML's configuration is its default network with the
\Krts{} group, and for v0\_2 it is its default network. The two
recombination rows compare v0\_1 and v0\_2 with the measurements they were
built from.}
\label{tab:same}
\end{table}

\section*{Where the versions differ without data}

At the notes' reference point, \lm{p}\unit{atm} of Kr with \lm{xF2}\,\%
F$_2$ pumped at \lm{P}\unit{kW\,cm^{-3}} for \lm{pulse}\unit{\mu s}, no
measurement exists, and there the defaults differ too. v0 sends
\rv{ref.v0}\,\% of the deposited energy out at 248\unit{nm}, v0\_1
\rv{ref.v01}\,\% and v0\_2 \rv{ref.v02}\,\%. Table~\ref{tab:ref} takes the
gaps apart. Between v0 and v0\_1 the default network accounts for
\rv{ref.d.network} of the \rv{ref.gap} points. v0 holds $n_e$ at
$10^{15}$\unit{cm^{-3}}, where v0\_1's free $n_e$ runs at
$\vm{neF2}$\unit{cm^{-3}} while F$_2$ lasts, so electrons quench
\lm{brf.e}\,\% of the KrF in v0 and \vm{brf.e}\,\% in v0\_1. v0 also leaves
\lm{out.ions}\,\% of the deposited energy in ions that never recombine.
v0\_1 returns that energy through recombination, mostly as heat, and its
gas ends at \vm{TgEnd}\unit{K} where v0's ends at \lm{TgEnd}\unit{K}. The
values v0\_1 changed account for the other \rv{ref.d.values} points. Only
the B-A branch moves the 248-nm share by more than a few hundredths, since
it sends \vm{pc.BA}\,\% of B's emission to 272\unit{nm}.

v0\_2 differs from v0\_1 by \rv{ref.gap12} points at 248\unit{nm}, but its
gas ends \rv{ref.cool12}\unit{K} cooler. After F$_2$ burns out, the \Krs{}
that v0\_1 turns into Penning ionization and heat forms \Krts{} in v0\_2
and leaves as 147-nm light, \rv{ref.v02.vuv}\,\% of the deposited energy.
With Kannari's two estimates v0\_2 gives \rv{ref.v02kannari}\,\%, what
v0\_1 gives with its optional \Krts{} group (\vi{kr2.on}\,\% in v0\_1's
overview).

\begin{table}[htbp]
\centering\footnotesize
\renewcommand{\arraystretch}{1.12}
\begin{tabular}{@{}lrrr@{}}
\toprule
model at the reference point & 248\unit{nm}, \% of \Edep & 147\unit{nm},
  \% of \Edep & $T_g$ at the end, K\\
\midrule
v0 by default, $n_e$ held at $10^{15}$\unit{cm^{-3}}, no charge chemistry &
  \rv{ref.v0} & \rv{ref.v0.vuv} & \rv{ref.v0.tg}\\
v0 with its ion and F$_2$-electron groups, $n_e$ free & \rv{ref.v0groups}
  & \rv{ref.v0groups.vuv} & \rv{ref.v0groups.tg}\\
\textbf{v0\_1} & \textbf{\rv{ref.v01}} & \textbf{\rv{ref.v01.vuv}} &
  \textbf{\rv{ref.v01.tg}}\\
v0\_1 without the B-A branch & \rv{ref.noBA} & \rv{ref.noBA.vuv} &
  \rv{ref.noBA.tg}\\
v0\_1 without F + F + Kr & \rv{ref.noFF} & \rv{ref.noFF.vuv} &
  \rv{ref.noFF.tg}\\
v0\_1 with v0's \Krp{} conversion and \Krtp{} recombination &
  \rv{ref.v0ions} & \rv{ref.v0ions.vuv} & \rv{ref.v0ions.tg}\\
v0\_1 without the ion-ion temperature law & \rv{ref.noIIT} &
  \rv{ref.noIIT.vuv} & \rv{ref.noIIT.tg}\\
v0\_1 with REBUILD's KrF and \KtF{} energies & \rv{ref.v0E} &
  \rv{ref.v0E.vuv} & \rv{ref.v0E.tg}\\
\textbf{v0\_2} & \textbf{\rv{ref.v02}} & \textbf{\rv{ref.v02.vuv}} &
  \textbf{\rv{ref.v02.tg}}\\
v0\_2 with Kannari's \Krts{} + \Krts{} & \rv{ref.v02pen} &
  \rv{ref.v02pen.vuv} & \rv{ref.v02pen.tg}\\
v0\_2 with Kannari's \Krts{} + F and \Krts{} + \Krts{} & \rv{ref.v02kannari}
  & \rv{ref.v02kannari.vuv} & \rv{ref.v02kannari.tg}\\
\bottomrule
\end{tabular}
\caption{The reference point, one change at a time. The rows below v0\_1
each undo one of its changes. The rows below v0\_2 give its two
unmeasured \Krts{} rates Kannari's estimates, and the last is v0\_1 with
its optional \Krts{} group.}
\label{tab:ref}
\end{table}

v0\_2 carries every rate that is known only as a published spread, or not
at all, as a band. A measured rate keeps its measurement as the default
and spans the other published values. A rate with no measurement is zero
by default and runs to the largest published value, which is no bound.
Table~\ref{tab:bands} runs each band alone at its two ends. Two of the
\rv{band.n} matter for the laser line, \Krts{} + F and the trimer law. The
\KtF{} quench bands raise the 248-nm share through the gas temperature,
since their energy goes to heat instead of 420-nm light and the hotter gas
slows the trimer.

\begin{table}[htbp]
\centering\footnotesize
\renewcommand{\arraystretch}{1.12}
\begin{tabularx}{\linewidth}{@{}>{\raggedright\arraybackslash}X
  >{\raggedright\arraybackslash}p{1.55cm}>{\raggedright\arraybackslash}p{5.2cm}rrr@{}}
\toprule
rate & v0\_2's value & band & 248\unit{nm} & 147\unit{nm} & $T_g$ end,
  K\\
\midrule
\rvBands
\bottomrule
\end{tabularx}
\caption{v0\_2's bands at the reference point. Each row runs one band at
its low end and then at its high end, the rest at their defaults, and
gives the percentages of \Edep{} and the final gas temperature at the two
ends. The unmeasured rates start at zero.}
\label{tab:bands}
\end{table}

Table~\ref{tab:grid} runs the three defaults over the notes' ranges, where
no measurement exists either. v0 and v0\_1 differ by up to
\rv{grid.diff}\,\% of v0\_1's value, and v0\_1 and v0\_2 by up to
\rv{grid.diff12}\,\%. \Krts{} + F at the top of its band raises v0\_2's
share by up to a factor \rv{grid.kr2}, and the trimer's weak law lowers it
by up to \rv{grid.tri}\,\%. Those two bands reach farther than the spread
between the versions at \rv{grid.nband} of the \rv{grid.n} points. The
spread is the wider only where v0\_2's gas ends at \rv{grid.hotmin}\unit{K}
or more, beyond every rate law's data. No version can be preferred there
on its numbers.

\begin{table}[htbp]
\centering\footnotesize
\renewcommand{\arraystretch}{1.08}
\begin{tabular}{@{}rrrrrrrr@{}}
\toprule
$p$ & $P$ & \multicolumn{5}{c}{248\unit{nm}, \% of \Edep} & $T_g$ at the
  end\\
\cmidrule(lr){3-7}
atm & kW\,cm$^{-3}$ & v0 & v0\_1 & v0\_2 & v0\_2, \Krts{} + F at
  $10^{-9}$ & v0\_2, trimer $T_g^{-1/2}$ & v0\_2, K\\
\midrule
\rvGrid
\bottomrule
\end{tabular}
\caption{The three defaults over the notes' pressure and power ranges,
0.5\,\% F$_2$ for 1\unit{\mu s}. The sixth and seventh columns put v0\_2's
two wide bands at the ends that move the share most.}
\label{tab:grid}
% the provenance note rides with the last table, so it ends the document
% wherever the floats land
\vspace{18pt}
{\normalsize\raggedright\color{ink2} Generated from LIGHT v0
(\lm{version}), v0\_1 (\vm{version}), v0\_2 (\rv{v02.version}) and ML by
\texttt{docs/}\allowbreak\texttt{make\_overview.py},
\texttt{docs/}\allowbreak\texttt{make\_overview\_v0\_1.py}, the two map
scripts and \texttt{docs/REVIEW/}\allowbreak\texttt{make\_review.py}. No
number in this document is typed by hand.\par}
\end{table}

\end{document}
